%! Summary Statistics for a Quantitative Variable — Unit 1, Lesson 1.5 — Slides (Beamer)
\documentclass[aspectratio=169, 11pt]{beamer}
\usepackage{apstats-beamer}   % provides \navyheader{} and \sectionlabel[color]{}
\usetikzlibrary{positioning}  % -beamer loads tikz but not this library

% The C and Y column types live in apstats-article, which a beamer deck does not
% load — redeclare them here rather than pulling in the article preamble.
\newcolumntype{C}[1]{>{\centering\arraybackslash}p{#1}}
\newcolumntype{Y}{>{\centering\arraybackslash}X}

\begin{document}

% TITLE SLIDE (CourseName is not defined in beamer, so the name is literal)
{
\setbeamercolor{background canvas}{bg=navy}
\begin{frame}[plain]
  \vspace{2.8cm}\hspace{0.55cm}%
  \begin{minipage}{0.9\textwidth}
    {\color{goldacc}\bfseries\small LESSON 1.5}\\[0.5cm]
    {\color{white}\bfseries\LARGE Summary Statistics for a Quantitative Variable}\\[1.2cm]
    {\color{skymid}\small Statistics~$\cdot$~Unit 1}
  \end{minipage}
\end{frame}
}

% GRADUAL RELEASE deck: title → targets → warm-up hook → I DO (notes frames) → WE DO (guided
% practice) → spoken debrief → close & START THE HOMEWORK, which is the individual
% practice — there is no independent practice set in the notes → AP Practice → close.
% There is NO group activity in this course — the notes carry the whole release. There is no
% transfer set and no reflection box: the water-use record lives only in the AP Practice page.
% Vocabulary is given up front here, not withheld: the notes frames ARE the direct instruction.
\newcommand{\redink}[1]{{\color{redacc}\bfseries #1}}

% ── LEARNING TARGETS ─────────────────────────────────────────────────────────
\begin{frame}[t, plain]
  \navyheader{Today's targets}
  \sectionlabel{Lesson 1.5}
  By the end of today, I can\ldots
  \begin{itemize}\setlength{\itemsep}{5pt}
    \item calculate the \textbf{mean} and the \textbf{median}, and find the \textbf{quartiles}
          by splitting an ordered list.
    \item put a single number on spread --- the \textbf{range}, the \textbf{IQR}, or the
          \textbf{standard deviation} --- and say which one unusual value can hijack.
    \item choose which pair of summaries to report and defend the choice from the
          \textbf{shape}, using the words \textbf{resistant} and \textbf{nonresistant}.
  \end{itemize}
  \begin{block}{How today works}
    Warm-up $\to$ \textbf{notes together} $\to$ \textbf{one we work as a class} $\to$
    debrief $\to$ \textbf{you start the homework here, alone}.
  \end{block}
\end{frame}

% ── HOOK ─────────────────────────────────────────────────────────────────────
{
\setbeamercolor{background canvas}{bg=navy}
\begin{frame}[plain]
  \vspace{1.5cm}
  \begin{center}
    {\color{skymid}\bfseries\small NINE PEOPLE. ONE PAYROLL. TWO ANSWERS.}\\[0.9cm]
    {\color{white}\bfseries\Large One summary says the typical worker here takes home
    \$1{,}000 a week.\\[0.4cm]
    The other says \$750.}\\[0.9cm]
    {\color{goldacc}\small Both are correct. Nobody in the building earns \$1{,}000.}
  \end{center}
\end{frame}
}

% ── I DO — direct instruction begins ─────────────────────────────────────────
{
\setbeamercolor{background canvas}{bg=navy}
\begin{frame}[plain]
  \vspace{1.8cm}\hspace{0.8cm}%
  \begin{minipage}{0.86\textwidth}
    {\color{goldacc}\bfseries\small GUIDED NOTES \ \textbullet\ \ I DO}\\[0.7cm]
    {\color{white}\bfseries\Large Open to your Guided Notes page. Fill each blank as we
    name it --- not before.}\\[0.9cm]
    {\color{skymid}\small Six terms go in the vocabulary box today. Write each one the
    moment we use it.}
  \end{minipage}
\end{frame}
}

% ── CENTRE: MEAN AND MEDIAN ──────────────────────────────────────────────────
\begin{frame}[t, plain]
  \navyheader{Two measures of center}
  \sectionlabel{notes, Center $\cdot$ a statistic is a numerical summary of sample data}
  \vspace{4pt}
  \begin{columns}[t]
    \begin{column}{0.47\textwidth}
      \begin{block}{Mean $\bar x$}
        Add every value; divide by the number of values (EK UNC-1.I.2).
      \end{block}
      \vspace{4pt}
      {\small Uses \textbf{all nine} paychecks.\\[3pt]
      \emph{\$9{,}000 $\div$ 9 $=$ \$1{,}000}}
    \end{column}
    \begin{column}{0.47\textwidth}
      \begin{block}{Median}
        The middle value of the \textbf{ordered} data; with an even count, average the two
        middle values (EK UNC-1.I.3).
      \end{block}
      \vspace{4pt}
      {\small Uses \textbf{one position}, not nine values.\\[3pt]
      \emph{the 5th paycheck $=$ \$750}}
    \end{column}
  \end{columns}
  \vspace{12pt}
  \begin{center}
    {\large Both are correct. \redink{Nobody earns \$1{,}000.}}\\[6pt]
    {\small\color{slate} Which one to \emph{report} is the question the Resistance row settles.}
  \end{center}
\end{frame}

% ── POSITION: QUARTILES AND PERCENTILES ─────────────────────────────────────
\begin{frame}[t, plain]
  \navyheader{Measures of position}
  \sectionlabel{notes, Quartiles $\cdot$ where does a value sit in the pack?}
  \vspace{4pt}
  \begin{itemize}\setlength{\itemsep}{7pt}
    \item Route A, twelve times: \; 27 \; 28 \; 28 \; 29 \; 29 \; 30 \; $|$ \; 30 \; 31 \; 31 \;
          32 \; 32 \; 33. \quad Median $= 30$.
    \item \textbf{$Q_1$} is the median of the lower half; \textbf{$Q_3$} is the median of the
          upper half (EK UNC-1.I.4). Here $Q_1 = 28.5$ and $Q_3 = 31.5$.
    \item Between them sits the \redink{middle 50\%} of the data.
    \item The \textbf{$p$th percentile} is the value with $p\%$ of the data \textbf{at or
          below} it (EK UNC-1.I.5). $Q_1$ and $Q_3$ are the 25th and 75th.
  \end{itemize}
  \vspace{8pt}
  \begin{block}{You already did this}
    Warm-up 3: \; 3 \; 5 \; 8 \; \textbf{9} \; 12 \; 15 \; 20. \;
    You circled 9, then 5 and 15. Those were the median, $Q_1$, and $Q_3$.
  \end{block}
\end{frame}

% ── SPREAD ───────────────────────────────────────────────────────────────────
\begin{frame}[t, plain]
  \navyheader{Three measures of variability}
  \sectionlabel{notes, Spread $\cdot$ spread, as a single number}
  \vspace{4pt}
  \begin{itemize}\setlength{\itemsep}{7pt}
    \item \textbf{Range} $=$ maximum $-$ minimum. One number, \redink{not} an interval
          (EK UNC-1.J.2). Route A: $33-27=6$ minutes --- not ``27 to 33.''
    \item \textbf{IQR} $= Q_3 - Q_1$ --- the width of the middle half. Route A: $31.5 - 28.5 = 3$.
    \item \textbf{Standard deviation} $s_x$ --- roughly the typical distance of a value from
          the mean (EK UNC-1.J.3). Its square $s_x^2$ is the \textbf{variance}.
  \end{itemize}
  \vspace{6pt}
  \begin{center}
    $\displaystyle s_x = \sqrt{\frac{1}{n-1}\sum_{i=1}^{n}(x_i-\bar x)^2}$
    \qquad\qquad
    {\small 26, 28, 30, 32, 34: \; $s_x^2 = \dfrac{40}{4} = 10$, \; $s_x \approx 3.16$}
  \end{center}
  \vspace{4pt}
  {\small\color{slate} We read $s_x$ off technology from here on. Change the units and
  \emph{every} one of these is multiplied by the same factor (EK UNC-1.J.4).}
\end{frame}

% ── THE CRUX, PART 1: ONE VALUE CHANGES ──────────────────────────────────────
\begin{frame}[t, plain]
  \navyheader{One value changes}
  \sectionlabel{notes, Resistance $\cdot$ the whole point of today}
  \vspace{6pt}
  \begin{columns}[t]
    \begin{column}{0.46\textwidth}
      \begin{block}{The owner's \$3{,}000 becomes \$750}
        Mean: \$1{,}000 $\to$ \redink{\$750}\\[3pt]
        Median: \$750 $\to$ \redink{\$750}
      \end{block}
    \end{column}
    \begin{column}{0.5\textwidth}
      \begin{block}{Route C is Route A with one late package}
        Range: 6 $\to$ \redink{24} --- the same as Route B\\[3pt]
        IQR: 3 $\to$ \redink{3}
      \end{block}
    \end{column}
  \end{columns}
  \vspace{14pt}
  \begin{center}
    {\large The mean adds every value. The median asks only \redink{which one is fifth.}}\\[6pt]
    {\small\color{slate} The range saw only the two ends. The IQR never touched the late
    package.}
  \end{center}
\end{frame}

% ── THE CRUX, PART 2: RESISTANT VS NONRESISTANT ──────────────────────────────
\begin{frame}[t, plain]
  \navyheader{Resistant or nonresistant --- and which pair to report}
  \sectionlabel{EK UNC-1.K.2 $\cdot$ Skill 4.B}
  \vspace{4pt}
  \begin{columns}[t]
    \begin{column}{0.47\textwidth}
      \begin{block}{Resistant}
        \textbf{Median} \quad\textbf{IQR}\\[3pt]
        One extreme value barely moves them.
      \end{block}
    \end{column}
    \begin{column}{0.47\textwidth}
      \begin{block}{Nonresistant}
        \textbf{Mean} \quad\textbf{Standard deviation} \quad\textbf{Range}\\[3pt]
        One extreme value drags them.
      \end{block}
    \end{column}
  \end{columns}
  \vspace{10pt}
  \begin{block}{The shape decides --- both measures of center are correct}
    Skewed, or an outlier present $\to$ report the \redink{median and IQR}. \quad
    Roughly symmetric, no outliers $\to$ report the \redink{mean and standard deviation}.
  \end{block}
  \vspace{4pt}
  {\small\color{slate} Two outlier rules (EK UNC-1.K.1): more than $1.5\times\text{IQR}$
  past a quartile, or 2 or more standard deviations from the mean. The first one goes to work
  next lesson.}
\end{frame}

% ── WE DO ────────────────────────────────────────────────────────────────────
\begin{frame}[t, plain]
  \navyheader{Guided Practice: The Quiz Scores}
  \sectionlabel[goldacc]{We do $\cdot$ you hold the pen}
  Eleven scores on a 20-point quiz: \; 12 \; 13 \; 14 \; 15 \; 15 \; 16 \; 17 \; 18 \; 18 \;
  19 \; 19.
  \begin{itemize}\setlength{\itemsep}{5pt}
    \item \textbf{14--15.} Fill the five-number summary: minimum, $Q_1$, median, $Q_3$, maximum;
          then the mean.
    \item \textbf{16--17.} The range and the IQR --- one number each; is the 12 an outlier?
    \item \textbf{18--20.} Then the lowest score turns out to be a \textbf{1}, not a 12. Which of
          \emph{mean, median, range, IQR} move? Which pair would you report now?
  \end{itemize}
  \begin{block}{The question behind the question}
    Which \redink{position} in the ordered list did we change --- and is that position anywhere
    near the middle?
  \end{block}
\end{frame}

% ── YOU DO — the homework is the individual practice, started in class ───────
\begin{frame}[t, plain]
  \navyheader{You do: start the homework}
  \sectionlabel{Last 8 minutes $\cdot$ on your own $\cdot$ finish it at home}
  \begin{itemize}\setlength{\itemsep}{7pt}
    \item Seven commutes: mean, median, and which one represents the office --- from the shape.
    \item Five values, one changed to 32: what moved and what did not.
    \item Twelve homework times: $Q_1$, $Q_3$, IQR --- and show the split.
    \item A percentile of birth weight, interpreted in one sentence.
  \end{itemize}
  \begin{block}{Silent and alone --- I am circulating}
    \emph{Item 2} is the one I am reading over your shoulder. If your median moved, flag me
    before you leave.
  \end{block}
\end{frame}

% ── DEBRIEF ──────────────────────────────────────────────────────────────────
\begin{frame}[t, plain]
  \navyheader{Debrief: ``the median is more accurate''}
  \sectionlabel{8 minutes $\cdot$ pens down, papers up --- we say these aloud}
  \vspace{6pt}
  \begin{columns}[t]
    \begin{column}{0.46\textwidth}
      \begin{block}{Almost right}
        The median is not more \emph{accurate} --- the mean of the corrected scores really
        is 15.
      \end{block}
    \end{column}
    \begin{column}{0.46\textwidth}
      \begin{block}{Right}
        The median is \redink{resistant}: it only asks which score is sixth. The mean
        touched the 1.
      \end{block}
    \end{column}
  \end{columns}
  \vspace{14pt}
  \begin{block}{Two things to say out loud}
    Name the person who earns the mean --- then the person who earns the median. Then give me
    a data set where the \redink{mean and standard deviation} are the right pair. Neither
    summary wins in general --- the shape decides.
  \end{block}
\end{frame}

% ── HOMEWORK (started in class, scored) ─────────────────────────────────────
\begin{frame}[t, plain]
  \navyheader{Homework --- start it now}
  \sectionlabel[goldacc]{5 exercises $\cdot$ scored $\cdot$ due the first class after two study halls}
  \begin{itemize}\setlength{\itemsep}{5pt}
    \item Seven commutes: mean, median, and which one represents the office.
    \item Five values, one changed to 32: what moved and what did not.
    \item Twelve homework times: $Q_1$, $Q_3$, IQR.
    \item A percentile of birth weight, interpreted in one sentence.
    \item One multiple-choice item, with a one-sentence justification.
  \end{itemize}
  \begin{block}{This one is graded}
    This is your \redink{scored} homework, on paper, due the first class after two study halls.
    Start with 1, 2, and 5 while I am still here to answer questions --- finish the rest at home.
  \end{block}
\end{frame}

% ── AP PRACTICE & PREVIEW (assigned, unscored) ──────────────────────────────
\begin{frame}[t, plain]
  \navyheader{AP Practice \& what's next}
  \sectionlabel[goldacc]{AP Practice --- assigned, not scored}
  \begin{itemize}\setlength{\itemsep}{5pt}
    \item Four multiple-choice items on 15 households' daily water use --- including which
          pair to report, and why.
    \item One free-response set: name the shape, then \emph{justify} which center and which
          spread you report.
    \item Optional: apply both outlier rules to the 620-gallon home and see whether they agree.
  \end{itemize}
  \begin{block}{Not graded --- but this is what the exam looks like}
    The AP Practice page is \textbf{not scored}. Work it for the reps, and bring what you
    get stuck on to study hall or office hours. The \redink{graded} homework is the section
    behind it.
  \end{block}
\end{frame}

% ── CLOSE ────────────────────────────────────────────────────────────────────
{
\setbeamercolor{background canvas}{bg=navy}
\begin{frame}[plain]
  \vspace{1.4cm}\hspace{0.8cm}%
  \begin{minipage}{0.86\textwidth}
    {\color{goldacc}\bfseries\small BEFORE YOU GO}\\[0.7cm]
    {\color{white}\bfseries\Large Mean and median are both correct. The shape of the
    distribution --- and whether one strange value is sitting in it --- is what decides which
    one you are allowed to quote.}\\[0.8cm]
    {\color{skymid}\small Next: Lesson 1.6 --- \emph{Boxplots and Comparing Distributions}.
    Those five numbers become a picture, and the $1.5\times\text{IQR}$ rule decides where its
    whiskers stop.}
  \end{minipage}
\end{frame}
}

\end{document}
